

\documentclass[11pt, leqno]{amsart}

\usepackage{amsmath}
\usepackage{amsthm}
\usepackage{amssymb}
\usepackage{mathrsfs}

\usepackage[colorlinks, linkcolor=black, citecolor=black, urlcolor=black, linktocpage=true]{hyperref}

%\usepackage{backref}


% THEOREM Environments-----------------------------------------
%
 \newtheorem{theorem}{Theorem}[section]

 \newtheorem{lemma}[theorem]{Lemma}
 \newtheorem{proposition}[theorem]{Proposition}

 \theoremstyle{remark}
 \newtheorem{remark}[theorem]{Remark}

\allowdisplaybreaks[3]
 
 \numberwithin{equation}{section}

\newcommand{\NN}{\ensuremath{\mathbb{N}}}
\newcommand{\ZZ}{\ensuremath{\mathbb{Z}}}
\newcommand{\RR}{\ensuremath{\mathbb{R}}}
\newcommand{\TT}{\ensuremath{\mathbb{T}}}
\newcommand{\XX}{\ensuremath{\mathbb{X}}}
\newcommand{\CC}{\ensuremath{\mathbb{C}}}
\newcommand{\Cont}{\ensuremath{\mathcal{C}}}
\newcommand{\PP}{\ensuremath{\mathcal{P}}}
\newcommand{\xx}{\ensuremath{\mathbf{x}}}
\newcommand{\BB}{\mathcal{B}}
\newcommand{\GG}{\mathcal{G}}
\newcommand{\LL}{\mathscr{L}}

\newcommand{\sgn}{\operatorname{sign}}
\newcommand{\supp}{\operatorname{supp}}
\newcommand{\spn}{\operatorname{span}}


% -----------------------
\title[Quasi-greedy bases  and Jacobi polynomials in $L_p$ spaces]{Unconditional and quasi-greedy bases in $L_p$ with applications to Jacobi polynomials Fourier series}
% -----------------------

% -----------------------
\author[F. Albiac]{Fernando Albiac}
\address{Mathematics Department\\ 
Universidad P\'ublica de Navarra\\
Campus de Arro\-sa\-d\'{i}a\\
Pamplona\\ 
31006 Spain}
% Tel.: +34-948-169553\\
% Fax: +34-948-166057\\
\email{fernando.albiac@unavarra.es}

\author[J. L. Ansorena]{Jos\'e L. Ansorena}
\address{Department of Mathematics and Computer Sciences\\
Universidad de La Rio\-ja\\ 
Logro\~no\\
26006 Spain}
% Tel.: +34-941299464 \\
% Fax: +34-941299460}
\email{joseluis.ansorena@unirioja.es}

\author[\'O. Ciaurri]{\'Oscar Ciaurri}
\address{Department of Mathematics and Computer Sciences\\
Universidad de La Rio\-ja\\ 
Logro\~no\\
26006 Spain}
% Tel.: +34-941299442 \\
% Fax: +34-941299460}
\email{oscar.ciaurri@unirioja.es}

\author[J. L. Varona]{Juan L. Varona}
\address{Department of Mathematics and Computer Sciences\\
Universidad de La Rio\-ja\\ 
Logro\~no\\
26006 Spain}
% Tel.: +34-941299451 \\
% Fax: +34-941299460}
\email{jvarona@unirioja.es}
% -----------------------

% -----------------------
\subjclass[2000]{46B15, 41A65}

\keywords{Thresholding greedy algorithm, unconditional basis, quasi-greedy basis, $L_p$-spaces, Jacobi polynomials}
% -----------------------

% -----------------------
\renewcommand{\datename}{\textbf{This paper has been published in}:}
\date{\textit{Rev.\@ Mat.\@ Iberoam.\@} \textbf{35} (2019), no.~2, 561--574.}
% -----------------------

% ------------------------------------------------------------------------
\begin{document}
% ------------------------------------------------------------------------

% -----------------------
\begin{abstract} 
We  show that  the decreasing rearrangement of the Fourier series with respect to  the Jacobi polynomials for functions in $L_p$  does not converge unless $p=2$. As a by-product of our work on quasi-greedy bases in $L_{p}(\mu)$, we show that no normalized unconditional basis in $L_p$, $p\not=2$, can be semi-normalized in $L_q$ for $q\not=p$, thus extending a classical theorem of Kadets and Pe{\l}czy{\'n}ski from 1962.
\end{abstract}
% -----------------------

% -----------------------
\maketitle
% -----------------------

% -----------------------
\section{Introduction and background}
\label{Introduction}
% -----------------------

A fundamental and total biorthogonal system for an infinite-dimensional  (real or complex) separable Banach space
$(\XX, \Vert\cdot\Vert)$
is a family $(\xx_{j},\xx_{j}^{\ast})_{j\in J}$ in  $\XX\times \XX^{\ast}$ verifying
\begin{itemize}
\item[(i)] $\XX= \overline{\spn \{ \xx_{j} : j\in J\}}$,
\item[(ii)] $\XX^*= \overline{\spn \{ \xx^*_{j} : j\in J\}}^{w^*}$,
 and
\item[(iii)] $\xx_{j}^{\ast}(\xx_{k})=1$ if $j=k$ and $\xx_{j}^{\ast}(\xx_{k})=0$ otherwise.
\end{itemize}
The family $\BB=(\xx_j)_{j\in J}$ is called a Markushevich basis (see~\cite{HMVZ2008}).
For brevity, we will refer to $\BB$ as a \textit{basis} and to the unequivocally determined collection of bounded linear functionals $\BB^*=(\xx_j^*)_{j\in J}$ as its
\textit{orthogonal family.}
If the basis verifies the additional condition
\begin{itemize}
\item[(iv)] $\sup_{j\in J} \Vert \xx_j\Vert \Vert \xx^*_j\Vert <\infty$
\end{itemize}
we say that  it is bounded in the sense of Markushevich, or $M$-\textit{bounded} for short. Finally, if we have
\begin{itemize}
\item[(v)] $0<\inf_{j\in J} \Vert \xx_j\Vert \le \sup_{j\in J} \Vert \xx_j\Vert <\infty$
\end{itemize}
we say that the basis 
is \textit{semi-normalized} (\textit{normalized} if $\Vert \xx_j\Vert=1$ for all $j\in J$). 
Notice that a basis verifies simultaneously (iv) and (v) if and only if
\[
\sup_{j \in J} \max\{ \Vert \xx_j\Vert ,\Vert \xx_j^*\Vert \}<\infty.
\]

Suppose $\BB=(\xx_j)_{j\in J}$ is a semi-normalized $M$-bounded basis in a Banach space $\XX$ with orthogonal family
$\BB^*=(\xx_j^*)_{j\in J}$.
 Each $f\in \XX$ is uniquely determined by its formal series expansion in terms of the basis,
\begin{equation}
\label{formula:SeriesExpansion}
f=\sum_{j\in J} \xx_j^*(f) \, \xx_j.
\end{equation}
In order to try to make sense of the infinite sum in \eqref{formula:SeriesExpansion}, one can fix a bijective mapping $\pi\colon\NN\to J$ and study the convergence of the formal series
$\sum_{n=1}^\infty \xx_{\pi(n)}^*(f)\xx_{\pi(n)}$. If this series converges to $f$ for every $f\in\XX$ then $\BB$ is a \textit{Schauder basis} for the bijection~$\pi$. Schauder  bases are very well-known and have been widely studied. They are characterized as those bases for which the partial sum operators $S_{\pi,m}\colon\XX\to\XX$, given by
\begin{equation}
\label{definition:PartialSum}
f\mapsto S_{\pi,m}(f)= \sum_{n=1}^m \xx_{\pi(n)}^*(f)\xx_{\pi(n)}
\end{equation}
are uniformly bounded. The property that  $\sum_{n=1}^\infty \xx_{\pi(n)}^*(f)\xx_{\pi(n)}$ converges for any $f\in\XX$ and any bijection $\pi$ yields the more restrictive
class of \textit{unconditional bases}. Recall that, equivalently,  a basis is unconditional if and only if
for every choice of signs $\varepsilon=(\varepsilon_j)_{j\in J}\in\{-1,1\}^J$ the multiplier
\[
P_\varepsilon\colon \XX\to \XX, \quad f\mapsto \sum_{j\in J} \varepsilon_j\xx_j^*(f) \xx_j
\]
is well defined and the family of operators $(P_\varepsilon)_{\varepsilon\in \{\pm 1\}^J}$ is uniformly bounded.

An \textit{ordering} for an element $f\in\XX$ (with respect to a basis $\BB$) is a one-to-one map $\rho\colon\NN\to J$ such that $\supp(f) := \{j\in J : \xx_j^*(f)\not=0\} \subseteq \rho(\NN)$.
From the point of view of  approximation theory, given a function $f$ in $\XX$ and an ordering $\rho$ for $f$, the sequence $(S_{\rho,m}(f))_{m=1}^\infty$
constructed as in \eqref{definition:PartialSum} defines an algorithm to approximate $f$. The minimal requirement we must impose to $\rho$ is that  $(S_{\rho,m}(f))_{m=1}^\infty$ converges to~$f$. In case $\BB$ is a Schauder basis for some bijection $\pi$, the algorithm based on $\pi$ fulfils this requirement for any $f\in\XX$. The independence of the ordering from the vector determines both the goodness and the limitations of this approximation algorithm for Schauder bases. The operators $S_{\pi,m}$ are linear and uniformly bounded,
but it is natural to wonder if by allowing the ordering to depend on each particular vector we can attain a higher rate of convergence.

The most important algorithm based on letting the ordering depend on the vector is the \textit{greedy algorithm}, also known as the
\textit{thresholding algorithm}. Since for each $f\in\XX$ the sequence $(\xx^*_j(f))_{j\in J}$ belongs to $c_0(J)$, there is an ordering $\rho$ for $f$ such that 
\begin{equation}
\label{greedycondition}
|\xx_{\rho(k)}^{\ast}(f)|\ge |\xx_{\rho(n)}^{\ast}(f)| \quad\text{if}\quad k\le n.
\end{equation}
If the family $(\xx_j^*(f))_{j\in J}$ contains several terms with the same absolute value then such an ordering for $f$ is not uniquely determined. 
In order to get uniqueness, we fix a ``natural'' bijection $\tau \colon J \to \NN$, and we impose the additional condition
\begin{equation}
\label{desempate}
\tau(\rho(k))\le\tau(\rho(n)) \quad\text{whenever}\quad |\xx_{\rho(k)}^{\ast}(f)|= |\xx_{\rho(n)}^{\ast}(f)|.
\end{equation}
If $f$ is infinitely supported, there is a unique ordering $\rho$ for $f$ that fulfils \eqref{greedycondition} and \eqref{desempate}, and such an ordering verifies $\rho(\NN)=\supp (f)$. In  case that $f$ is finitely supported, there is a unique ordering $\rho$ for $f$ that verifies \eqref{greedycondition}, \eqref{desempate}, and the extra property $\rho(\NN)=J$. In any case, we will refer to such a unique ordering as the \textit{greedy ordering} for~$f$.
For each $m\in \NN$, the \textit{$m$-term greedy approximation} to $f$
is given by
\[
  \GG_{m}[\BB, \XX](f):=\GG_{m}(f) = S_{\rho,m}(f) = \sum_{n=1}^{m}\xx_{\rho(n)}^{\ast}(f)\xx_{\rho(n)},
\]
where $\rho$ is the greedy ordering for $f$, and the sequence $(\GG_m(f))_{m=1}^\infty$ is called the greedy algorithm for $f$ with respect to the basis~$\BB$.

Konyagin and Temlyakov \cite{KonyaginTemlyakov1999} defined a basis $\BB$ to be \textit{quasi-greedy}
  if for $f\in \XX$,
$\lim_{m\to \infty}\GG_{m}(f)=f$, 
that is, the greedy algorithm with respect to  $\BB$ converges in the Banach space~$\XX$. Subsequently, Wojtaszczyk \cite{Wo2000} proved that these are precisely the bases for which the greedy operators $(\GG_m)_{m=1}^{\infty}$ are uniformly bounded, i.e., there exists a constant $C\ge 1$ such that, for all $f\in \XX$ and $m\in \NN$,
\begin{equation}
\label{result:WoCondition}
\Vert \GG_m(f)\Vert\le C \Vert f\Vert.
\end{equation}
Notice the similarity between \eqref{result:WoCondition} and the uniform boundedness  of the partial sum projections that characterizes Schauder bases.
However,  the operators $\GG_m$ are neither linear nor continuous. 
We emphasize that, as Wojtaszczyk pointed out in \cite{Wo2000}, the choice of the bijection $\tau$ with respect to which we construct the greedy algorithm $(\GG_{m})_{m=1}^{\infty}$ plays no relevant role in the theory.

Unconditional bases are a special kind of quasi-greedy bases. Although the converse is not true in general, quasi-greedy bases always retain a flavor of unconditionality,  in a certain sense. For instance, they are \textit{unconditional for constant coefficients} \cite{Wo2000}, i.e., there is a constant $C$  such that
\begin{equation}
\label{result:UCCCondition}
 \frac{1}{C}\left\Vert \sum_{j\in A} \xx_{i}\right\Vert\le \left\Vert \sum_{j\in A}\varepsilon_{i} \xx_{i}\right\Vert\le C \left\Vert \sum_{j\in A} \xx_{i}\right\Vert
\end{equation}
for any finite subset $A$ of $J$ and any choice of signs $\varepsilon_j\in\{\pm1\}$. To be precise, the constant given by $C=2 C_{w}$ works, where $C_{w}$ is the least constant in \eqref{result:WoCondition}.

Before the concept of quasi-greedy basis was introduced in the literature, C\'or\-do\-ba and Fer\-n\'an\-dez \cite{CoFe1998} had studied the convergence of decreasing rearranged Fourier series. For $k\in\ZZ$ let us define $\tau_k\colon\RR\to\CC$ by $\tau_k(x)=e^{2\pi k x i}$. Let $1\le p<\infty$ and denote by $q$ its conjugate exponent, determined by
$1/p+1/q=1$. Then, with the usual identification of  $L_p^*(\TT)$ with $L_q(\TT)$, the double sequence $(\tau_k,\tau_{-k})_{k=-\infty}^\infty$ is a semi-normalized $M$-bounded biorthogonal system for $L_p(\TT)$. The authors of \cite{CoFe1998} showed that for each $1<p<2$ there is a function $f\in L_p(\TT)$ whose decreasing rearranged Fourier series does not converge, which in our language can be stated as saying that  the trigonometric system $(\tau_k)_{k=-\infty}^\infty$ is not a quasi-greedy basis for $L_p(\TT)$.
Combining the condition characterizing quasi-greedy bases \eqref{result:WoCondition} with Remark~2 in \cite{Telmyakov1998}, the result extends to the whole range of $p\in[1,\infty]\setminus\{2\}$ (for $p=\infty$, replace $L_p(\TT)$ with $\Cont(\TT)$).
Wojtaszczyk gave a different proof of this result in \cite{Wo2000} that relies on~\eqref{result:UCCCondition}.

A natural way to continue this line of research is to consider Fourier series with respect to orthonormal bases. Let $(X,\Sigma,\mu)$ be a measure space such that the Hilbert space $L_2(\mu)$ is separable. Let
$(\xx_j)_{j\in J}$ be an orthonormal basis of $L_2(\mu)$.
For $1\le p< \infty$, let $q$ be its conjugate exponent, so that 
$L_p^*(\mu)$ is identified with $L_q(\mu)$.
In case that  $\spn \{ \xx_j : j\in J\}$ is dense in $L_p(\mu)$ and 
 weak${}^*$ dense in $L_q(\mu)$
we have that $(\xx_j)_{j\in J}$ is a basis for $L_p(\mu)$ with orthogonal family 
$(\overline{\xx_j})_{j\in J}$.
 It therefore makes sense to investigate the convergence of the greedy algorithm with respect to the Fourier series of $(\xx_j)_{j\in J}$.
Semi-normalized, $M$-bounded bases provide the natural framework to study quasi-greedy bases (see~\cite{Wo2000}).
Thus, an early test to discard a  basis $(\xx_j)_{j\in J}$ from being quasi-greedy (even under a suitable re-scaling) is to check if it is $M$-bounded, i.e., if it  fulfils the condition
\begin{equation}
\label{EarlyTest}
\sup_{j\in J} \Vert \xx_j \Vert_p \Vert \xx_j \Vert_q<\infty.
\end{equation}
Since  $(\xx_j)_{j\in J}$ need not to be semi-normalized in $L_p(\mu)$, 
it is natural to investigate the greedy algorithm with respect to the $L_p(\mu)$-normalized basis 
$(\Vert \xx_j\Vert_p^{-1} \xx_j)_{j\in J}$. This is how the greedy algorithm of the Haar system in $L_p$ is studied, for example.
 

Notice that if the measure $\mu$ is finite and the orthonormal basis $(\xx_j)_{j\in J}$
is \textit{uniformly bounded}, i.e., $\sup_{j\in J} \Vert \xx_j\Vert_\infty<\infty$, then it is semi-normalized and $M$-bounded in $L_p(\mu)$ for any $1\le p<\infty$.
Nielsen \cite{Nielsen2007} proved that there is an uniformly bounded orthonormal basis of $L_2(\TT)$ which is quasi-greedy for $L_p(\TT)$ for any $1<p<\infty$, thus
exhibiting a behavior opposite to that of the trigonometric system.

In this paper we focus on Jacobi polynomials.
Recall that, for  scalars $\alpha$, $\beta>-1$, the  $L_2(\mu_{\alpha,\beta})$-normalized Jacobi polynomials
$(p_n^{(\alpha,\beta)})_{n=0}^\infty$ appear as the orthonormal polynomials associated to the measure $\mu_{\alpha,\beta}$ given by
\begin{equation}
\label{measureJacobi}
d\mu_{\alpha,\beta}(x)=(1-x)^\alpha (1+x)^\beta \chi_{(-1,1)} (x) \, dx.
\end{equation}
Since polynomials are dense in $L_p(\mu_{\alpha,\beta})$ for any $1\le p<\infty$, Jacobi polynomials of indices $\alpha$ and $\beta$ constitute an orthonormal basis of $L_2(\mu_{\alpha,\beta})$. Our main result on Jacobi polynomials establishes that the greedy algorithm for this kind of orthogonal polynomials follows the same pattern as the greedy algorithm for the trigonometric system.

\begin{theorem}
\label{result:MainTheorem}
Let $1\le p<\infty$ and $\min\{\alpha,\beta\}> -1/2$. The $L_p(\mu_{\alpha,\beta})$-normalized Jacobi polynomials of indices $\alpha$ and $\beta$ form a quasi-greedy basis for $L_p(\mu_{\alpha,\beta})$ if and only if $p=2$.
\end{theorem}


Section~\ref{Jacobi} is devoted to proving Theorem~\ref{result:MainTheorem}.
Before, in Section~\ref{UnconditionalLp}, we develop
the functional analysis machinery that we will need in order to do that and 
we show the following result on unconditional bases in $L_p(\mu)$-spaces.

\begin{theorem}
\label{NotSimultaneouslyBounded}Let $\mu$ be a finite measure and $p\in(1,\infty)\setminus\{2\}$. Suppose that $(\xx_j)_{j\in J}$ is a semi-normalized unconditional basis of a non-Hilbertian Banach space $\XX\subseteq L_p(\mu)$. Suppose also that $\XX$ is complemented in $L_p(\mu)$. Then
\begin{itemize}
\item[(i)] $\limsup_{j\in J} \Vert \xx_j\Vert_q=\infty$ for any $p<q$, and
\item[(ii)] $\liminf_{j\in J} \Vert \xx_j\Vert_q=0$ whenever $\max\{q,2\}<p$.
\end{itemize}
\end{theorem}

Notice that, since the basis need not  be unconditional in the 
$L_q$-norm, part (ii) in the above theorem cannot be deduced from part (i). We emphasize that Theorem~\ref{NotSimultaneouslyBounded} is relevant for its own intrinsic importance within the theory of bases. Firstly, it extends to any $q$ a result that Kadets and Pe{\l}czy{\'n}ski proved only for $q=2$ (see Corollary~9 in  \cite{KadPel1962}). Secondly, it  generalizes the main result of Gapo\v{s}kin in \cite{Gaposkin1958}, where he shows that no normalized unconditional basis in
$L_p[0,1]$ can be uniformly bounded. Lastly, for finite measures, Theorem~\ref{NotSimultaneouslyBounded} overrides a recent result of the first two authors that says that if $\mu$ is a nonpurely atomic measure then there is no basis $\BB$ that is simultaneously greedy (see the definition below) in two different $L_{p}(\mu)$ spaces, $1<p<\infty$ (see Theorem~4.4 in \cite{AA2015}). 


We end this preliminary section by singling out some notation and terminology that will be used heavily throughout.
Given families of positive real numbers $(\alpha_i)_{i\in I}$ and $(\beta_i)_{i\in I}$, the symbol $\alpha_i\lesssim \beta_i$ for $i\in I$ means that
$\sup_{i\in I}\alpha_i/\beta_i <\infty$, while $\alpha_i\approx \beta_i$ for $i\in I$ means that $\alpha_i\lesssim \beta_i$ and $\beta_i\lesssim \alpha_i$ for $i\in I$.

A basis $\BB=(\xx_j)_{j\in J}$ in a Banach space $\XX$ is said to be \textit{democratic} if
there is a constant $D\ge 1$ such that
\[
\left\Vert \sum_{j\in A}\xx_j \right\Vert \le D \left\Vert \sum_{j\in B} \xx_j \right\Vert
\]
whenever $A$ and $B$ are finite subsets of $J$ with $|A|=|B|$.
To quantify the democracy of a basis $\BB$ we consider the \textit{upper democracy function} of~$\BB$ (also known as the \textit{fundamental function of~$\BB$}) given by
\[
 \varphi_{u}[\BB,\XX](N)=\sup_{|A|\le N}\left\Vert \sum_{j\in A}\xx_j\right\Vert, \qquad N\in\NN,
\]
and the \textit{lower democracy function} of $\mathcal B$ in $\XX$, defined as
\[
 \varphi_{l}[\BB,\XX](N)=\inf_{|A|\ge N}\left\Vert \sum_{j\in A}\xx_j\right\Vert, \qquad N\in\NN.
\]
A quasi-greedy basis $\mathcal B$ is democratic if and only if $\varphi_{u}[\BB,\XX](N)\approx \varphi_{l}[\BB,\XX](N)$ for $N\in\NN$.

A basis $(\xx_n)_{n=1}^\infty$ is said to be \textit{almost greedy} if there is a constant $C\ge 1$ such that for all $m\in \NN$ and all $x\in \XX$ we have
\[
 \Vert x-\GG_m(x)\Vert \le C \inf\left\{\left \Vert x-\sum_{j\in A} \xx_j^*(x) \xx_j\right\Vert : |A|=m\right\}.
\]
Dilworth et al.\ \cite{DKKT2003} characterized almost greedy basis as those bases that are simultaneously quasi-greedy and democratic.

Finally, the best one can hope for in regards to the greedy algorithm is the existence of a constant $C\ge 1$ such that for all $m\in \NN$ and all $x\in \XX$,
\[
 \Vert x-\GG_m(x)\Vert \le C \inf\left\{\left \Vert x-\sum_{j\in A} a_{j} \xx_j\right\Vert : |A|=m,\; (a_{j})_{j\in A}\;\;\text{scalars}\right\}.
 \]
If this is the case, the basis is called \textit{greedy}. Konyagin and Temlyakov \cite{KonyaginTemlyakov1999} characterized greedy bases as those bases that are unconditional and democratic.

If necessary, the reader will find more background on Banach space theory and greedy bases in \cite{AlbiacKalton2006}, 
and on orthogonal polynomials in~\cite{Szego}. 

% -----------------------
\section{Quasi-greedy and unconditional bases in $L_p(\mu)$-spaces}
\label{UnconditionalLp}
% -----------------------

We get started by generalizing to quasi-greedy bases a fact that is standard for unconditional bases in $L_p(\mu)$-spaces.

\begin{lemma}
\label{result:Random}
Let $(X,\Sigma,\mu)$ be a finite measure space. Let $1\le p<\infty$ and $(\xx_j)_{j\in J}$ be a quasi-greedy basis for
a separable subspace of $L_p(\mu)$. Then for $A\subseteq J$ finite,
\[
\left\Vert \sum_{j\in A}\xx_j \right\Vert_p
\approx \left\Vert \left(\sum_{j\in A} |\xx_j|^2\right)^{1/2}\right\Vert_p.
\]
\end{lemma}

\begin{proof}
Let $(\varepsilon_j)_{j\in J}$ be a Rademacher family defined on some probability space $(\Omega,P)$, and let $A\subseteq J$ finite. Combining \eqref{result:UCCCondition}, Fubini's
theorem, and Khintchine's inequalities yields 
\begin{align*}
\left\Vert \sum_{j\in A}\xx_j \right\Vert_p
&\approx \left( \int_\Omega \Big\Vert \sum_{j\in A}\varepsilon_j \xx_j \Big\Vert_p^p\, dP\right)^{1/p}\\
&= \left(\int_X \int_\Omega \Big| \sum_{j\in A} \varepsilon_j \xx_j\Big|^p \, dP \,d\mu\right)^{1/p}\\
&\approx \left\Vert \left(\sum_{j\in A} |\xx_j|^2\right)^{1/2}\right\Vert_p.
\qedhere
\end{align*}
\end{proof}


Our next auxiliary result displays an estimate that is implied when a family of functions is simultaneously seminormalized in two different $L_{p}$ spaces.

\begin{lemma}
\label{result:PreviousDemocracyEstimates} 
Let $1\le p< q\le 2$ (respectively, $2\le q <p\le\infty$) and let $(f_j)_{j\in J}$ be a family of measurable functions defined on a finite measure space $(X,\Sigma,\mu)$. Suppose that
$\Vert f_j\Vert_p\approx \Vert f_j \Vert_q\approx 1$ for $j\in J$. Then, for $A\subseteq J$ finite,
\[
|A|^{1/2} \lesssim \left\Vert \left( \sum_{j\in A}| f_j|^2\right)^{1/2}\right\Vert_p \lesssim |A|^{1/q}
\]
(respectively,
\[
|A|^{1/q} \lesssim \left\Vert \left( \sum_{j\in A}| f_j|^2\right)^{1/2}\right\Vert_p \lesssim |A|^{1/2}).
\]
\end{lemma}

\begin{proof}
Assume $1\le p<q\le 2$. Using the embeddings $\ell_q\subseteq\ell_2$ and $L_q(\mu)\subseteq L_p(\mu)$,
\begin{equation*}
\left\Vert \left(\sum_{j\in A} | f_j|^2\right)^{1/2}\right\Vert_p
\lesssim \left\Vert \left(\sum_{j\in A} | f_j|^q\right)^{1/q}\right\Vert_q
=\left(\sum_{j\in A} \Vert f_j\Vert_q^q\right)^{1/q}
\approx |A|^{1/q}.
\end{equation*}
Let $r=p/2<1$. Using that $\Vert f+g\Vert_r \ge \Vert f\Vert_r +\Vert g\Vert_r$ whenever $f$ and $g$ are positive measurable  functions,
\begin{align*}
 \left\Vert \left(\sum_{j\in A} | f_j|^2\right)^{1/2}\right\Vert_p
&= \left\Vert \sum_{j\in A} | f_j|^2\right\Vert_r^{1/2}\\
&\ge \left( \sum_{j\in A} \Vert \,  |f_j|^2 \, \Vert_r \right)^{1/2}\\
&=\left( \sum_{j\in A} \Vert f_j \Vert_p^2 \right)^{1/2}\\
&\approx |A|^{1/2}.
\end{align*}

The case $2\le q<p\le\infty$ follows from a ``dual'' argument.
\end{proof}

\begin{lemma}
\label{result:DemocracyEstimates}
Let $(X,\Sigma,\mu)$ be a finite measure space. Suppose $1\le p< q\le 2$ (respectively, $2\le q <p<\infty$). Let  $(\xx_j)_{j\in J}$ be a quasi-greedy basis for
a separable subspace $\XX$ of $L_p(\mu)$ such that $\Vert \xx_j \Vert_q\approx 1$ for $j\in J$. Then, for $N\in\NN$,
\[
N^{1/2} \lesssim \varphi_{l}[\BB,\XX](N) \le \varphi_{u}[\BB,\XX](N) \lesssim N^{1/q}
\]
(respectively,
\[
N^{1/q} \lesssim \varphi_{l}[\BB,\XX](N) \le \varphi_{u}[\BB,\XX](N) \lesssim N^{1/2}).
\]
\end{lemma}

\begin{proof} 
Quasi-greedy bases are semi-normalized, i.e., $\Vert \xx_j \Vert_p\approx 1$ for $j \in J$. 
Now we just need to put together Lemma~\ref{result:Random} and Lemma~\ref{result:PreviousDemocracyEstimates}.
\end{proof}


The next two propositions are on-the-spot corollaries of Lemmas~\ref{result:PreviousDemocracyEstimates} and \ref{result:DemocracyEstimates}, respectively. We point out that a similar statement to Proposition~\ref{result:Democracy} with the stronger assumption that the basis be uniformly bounded was obtained by Dilworth et al.\ in Proposition~2.17 of~\cite{DST2012}.

\begin{proposition}
\label{result:LdosSumEstimate} 
Let $1\le p\le \infty$. Suppose $(f_j)_{j\in J}$ is a family of measurable functions defined on a finite measure space $(X,\Sigma,\mu)$ such that
$\Vert f_j\Vert_p\approx \Vert f_j \Vert_2\approx 1$ for $j\in J$. Then for $A\subseteq J$ finite,
\[
 \left\Vert \left( \sum_{j\in A}| f_j|^2\right)^{1/2}\right\Vert_p \approx |A|^{1/2}.
\]
\end{proposition}

\begin{proposition}
\label{result:Democracy}
Let $(X,\Sigma,\mu)$  be a finite measure space and let $1\le p<\infty$. Suppose  $\BB=(\xx_j)_{j\in J}$ is a quasi-greedy basis for
a separable subspace $\XX$ of $L_p(\mu)$ with $\Vert \xx_j \Vert_2\approx 1$ for $j\in J$. 
Then $\BB$ is democratic, hence almost greedy, and
its democracy functions verify 
\[
\varphi_l[\BB,\XX](N)\approx N^{1/2}\approx \varphi_u[\BB,\XX](N),\quad N\in\NN.
\]
\end{proposition}

We are now en route to completing the proof of Theorem~\ref{NotSimultaneouslyBounded}. Before we do so, we write down two classical results in the isomorphic theory of Banach spaces that are  well known to the specialists. In order to make the paper as self-contained as possible we sketch their proofs.

\begin{theorem}
\label{result:KaPeSubbasis} 
Let $(X,\Sigma,\mu)$ be a measure space. Suppose that $\BB$ is a seminormalized unconditional basis of a non-Hilbertian Banach space $\XX\subseteq L_p(\mu)$, $1<p<\infty$. Suppose also that $\XX$ is complemented in $L_p(\mu)$. Then
$\BB$ has a subbasis that is equivalent to the unit vector basis of~$\ell_p$.
\end{theorem}

\begin{proof}
 Without loss of generality we may and do assume that $L_p(\mu)$ is separable. Then $L_p(\mu)$ is isomorphic either to $L_p[0,1]$ or to~$\ell_p$
(see~\cite{JohnsonLindes2001}).
In the former case, the same argument used by Kadec and Pe{\l}czy{\'n}ski to prove Theorem~4 of \cite{KadPel1962} leads to our goal. In the latter, by the Bessaga-Pe{\l}czy{\'n}ski selection principle (see p.~214 of \cite{BePe1958}), $\BB$ has a subbasis equivalent to a block basic sequence of the unit vector basis of~$\ell_p$.
Since the unit vector basis of $\ell_p$ is perfectly homogeneous, this subbasis is equivalent to the unit vector basis of~$\ell_p$.
\end{proof}


\begin{lemma}
\label{result:EquivalentNorms} 
Let $(X,\Sigma,\mu)$ be a finite measure space and let $0<p<q\le \infty$. Consider a subset $S\subseteq L_p(\mu)$ such that $\Vert f\Vert_p\approx\Vert f\Vert_q$ for $f\in S$. Then, for any $0<r\le q$ we have $\Vert f\Vert_r \approx\Vert f\Vert_p\approx\Vert f\Vert_q$ for $f$ in~$S$.
\end{lemma}

\begin{proof} 
The result is obvious for $p\le r\le q$, so we assume that $r<p$. Then, it is also obvious that $\Vert f\Vert_r\lesssim \Vert f\Vert_p$ for $f\in S$. To prove the reverse inequality, consider $0<a<p$ such that $a/q+(p-a)/r=1$.
 For $f\in S$, by H\"older's inequality,
\begin{align*}
\Vert f\Vert_p^p 
&=\int_X |f|^a |f|^{p-a} \, d \mu\\
&\le \left( \int_X |f|^{q}\, d\mu \right)^{a/q} \left( \int_X |f|^{r}\, d\mu \right)^{(p-a)/r}\\
&=\Vert f\Vert_q^a \Vert f\Vert_r^{p-a}\\
&\lesssim \Vert f\Vert_p^a \Vert f\Vert_r^{p-a}.
\end{align*}
After simplifying we get $\Vert f\Vert_p^{p-a}\lesssim \Vert f\Vert_r^{p-a}$.
\end{proof}


\begin{proof}[Proof of Theorem~\ref{NotSimultaneouslyBounded}] 
Let $(\xx_j)_{j\in J}$ be a semi-normalized unconditional basis of a non-Hilbertian Banach space $\XX\subseteq L_p(\mu)$, where $p\in(1,\infty)\setminus\{2\}$ and $\mu$ is finite.
We divide the proof in three cases.

\smallskip

\noindent\textsc{Case 1:} $1<p<2$ and $p<q$. Assume that $\limsup_j \Vert x_j\Vert_q<\infty$.
Without loss of generality we can suppose that $p<q\le 2$.
Put $J_0=\{ j : \Vert \xx_j\Vert_q<\infty \}$, let $\BB_0=(\xx_j)_{j\in J_0}$, and then define
$\XX_0$ as the closed subspace spanned by $\BB_0$ in $L_p(\mu)$.
We have that $J\setminus J_0$ is finite and that
$\Vert \xx_j\Vert_q\approx 1$ for $j\in J_0$. By Lemma~\ref{result:DemocracyEstimates}, $\varphi_u[\BB_0,\XX_0](N)\lesssim N^{1/q}$. Furthermore, by
Theorem~\ref{result:KaPeSubbasis}, $\BB_0$ has a subbasis equivalent to the unit vector basis of $\ell_p$. Therefore,
$N^{1/p}\lesssim \varphi_u[\BB_0,\XX_0](N)$. Combining, we obtain $N^{1/p}\lesssim N^{1/q}$. This absurdity proves the result. %for $1<p<2$.

\smallskip

\noindent\textsc{Case 2:} $2<p<\infty$ and $q<p$. This case is the ``dual'' of the previous one. Since  its proof is similar we leave it for the reader.

\smallskip

\noindent\textsc{Case 3:} $2<p<q$. Suppose that $\limsup_j \Vert x_j\Vert_q<\infty$. Removing a finite set of terms from $\BB$ we get
$\Vert \xx_j\Vert_q\approx 1$. By Lemma~\ref{result:EquivalentNorms}, $\Vert \xx_j\Vert_2 \approx 1$, contradicting the already proven Case~2.
\end{proof}

\begin{remark} 
The proof of Theorem~\ref{NotSimultaneouslyBounded}
reinforces the role of  democracy as a hinge property in the study of unconditional bases in Banach spaces. This idea was already inferred from the work of Zippin~\cite{Zippin1966}, where he characterized perfectly homogeneous bases.
\end{remark}


\begin{remark} 
The subspace spanned by the Rademacher functions in $L_{p}(\mu)$ serves as an example to show that the assumption
``non-Hilbertian'' cannot be dropped from
Theorem~\ref{NotSimultaneouslyBounded}.
\end{remark}

% -----------------------
\section{The greedy algorithm for Jacobi polynomials}
\label{Jacobi}
% -----------------------

In this section, besides the orthonormal polynomials $(p_n^{(\alpha,\beta)})_{n=0}^\infty$ defined in Section~\ref{Introduction}, we consider
the polynomials $(P_n^{(\alpha,\beta)})_{n=0}^\infty$, which are orthogonal for the measure 
$\mu_{\alpha,\beta}$ defined in \eqref{measureJacobi} and verify the \textit{normalization condition}
\begin{equation}
\label{ValueInOne}
P_n^{(\alpha,\beta)} (1)=\binom{n+\alpha}{n}, \quad n\in \NN\cup\{0\}.
\end{equation}
Of course, there are positive scalars $d_n$ such that $p_n^{(\alpha,\beta)}=d_n P_n^{(\alpha,\beta)}$, $n\ge 0$. It is well known that the normalization sequence 
$(d_n)_{n=0}^\infty $ verifies
\begin{equation}
\label{eq:asym-dn}
d_n
= \Big(\frac {(2n+\alpha+\beta+1)n!\,\Gamma(n+\alpha+\beta+1)}
{2^{\alpha+\beta+1}\Gamma(n+\alpha+1)\Gamma(n+\beta+1)}
 \Big)^{1/2}
\approx n^{1/2}, \ n\in\NN.
\end{equation}
In what follows, with the aim to avoid cumbrous notations, we will denote by $(n^{1/2} P_n^{(\alpha,\beta)})_{n=0}^\infty$ the extension of $(n^{1/2} P_n^{(\alpha,\beta)})_{n=1}^\infty$ whose $0$-term
is the constant function~$1$.
In the light of \eqref{eq:asym-dn},
it is reasonable to expect that the sequences $(p_n^{(\alpha,\beta)})_{n=0}^\infty$ and $(n^{1/2} P_n^{(\alpha,\beta)})_{n=0}^\infty$ behave similarly.
Wojtaszczyk \cite{Wo2000} confirmed this fact by showing that quasi-greedy bases verify the following perturbation principle. 

\begin{theorem}[{[Proposition~3 in \cite{Wo2000}}]
\label{result:Perturbing} 
Suppose  $(\xx_j)_{j\in J}$ is a quasi-greedy basis for a Banach space~$\XX$. Let $(\lambda_j)_{j\in J}$ be a family of scalars such that
\[
0<\inf_{j\in J} |\lambda_j|\le \sup_{j\in J} |\lambda_j|<\infty.
\]
Then $(\lambda_j \xx_j)_{j\in J}$ is a quasi-greedy basis for $\XX$.
\end{theorem}


A powerful tool to carry out estimates involving Jacobi polynomials is the so-called \textit{Darboux formula}. The next theorem establishes an expression for the 
error term associated to this formula that is accurate enough for our purposes.

\begin{theorem}[Theorem~8.21.13 in \cite{Szego}]
\label{result:Darboux}
Let $\alpha$, $\beta>-1$ and $\delta>0$. Then
\begin{equation*}
n^{1/2} P_n^{(\alpha,\beta)}(\cos \theta)
= k(\theta)  \cos(n\theta+\phi(\theta))+ E_n(\theta) 
\end{equation*}
with 
\begin{equation*}
\begin{aligned}
\phi(\theta)
&= (\alpha+\beta+1)\theta/2-(2 \alpha+1)\pi/4, \\
k(\theta)
&= \pi^{-1/2}(\sin \tfrac{\theta}{2})^{-\alpha-1/2}
(\cos \tfrac{\theta}{2})^{-\beta-1/2},
\end{aligned}
\end{equation*}
and the error term $E_n(\theta)$ verifies
\[
E_n(\theta) = \frac{k(\theta)}{n \sin \theta} \, O(1)
\]
for $n\in\NN$, where the $O(1)$ holds uniformly in the interval 
${\delta}/{n}\le \theta \le \pi-{\delta}/{n}$.
\end{theorem}


Darboux formula provides tight estimates for Jacobi polynomials when the variable is not too close to the endpoints  $-1$ and $1$. The technique to estimate Jacobi polynomials near $1$ is also well-known for experts. It is based on the formula
\begin{equation}
\label{derivative}
\left(P_n^{(\alpha,\beta)}\right)'= (1+\alpha+\beta+n) P_{n-1}^{(\alpha+1,\beta+1)}
\end{equation}
and the behavior of the roots of Jacobi polynomials. In the following lemma we reproduce this standard argument for the sake of completeness.

\begin{lemma}
\label{result:NearOneEstimate}
Let $\alpha>-1$ and $\beta>-1$. There is $d>0$ such that
\[
P_n^{(\alpha,\beta)}(x)\approx n^\alpha
\]
for $n\in\NN$ and $1-d/n^2 \le x \le 1$.
\end{lemma}

\begin{proof}
Let  $z_n$ denote the largest root of $ P_n^{(\alpha,\beta)}$ and let $\gamma_n\in(0,\pi)$ be such that $\cos(\gamma_n)=z_n$.
It is known (see Theorem~8.9.1 in \cite{Szego}) that $\gamma_n\approx 1/n$, hence $1-z_n\approx 1/n^2$. 
Moreover, it is easy to deduce from \eqref{ValueInOne} and \eqref{derivative} that 
\[
P_n^{(\alpha,\beta)}(1)\approx n^\alpha \quad\text{and}\quad 
(P_n^{(\alpha,\beta)})'(1)\approx n^{\alpha+2}.
\] 
Choosing $d>0$ small enough we get
$z_n \le 1-d/n^2$ and
\[
n^\alpha\lesssim P_n^{(\alpha,\beta)}(1) - \frac{d}{n^2} (P_n^{(\alpha,\beta)})'(1).
\]
Let $1-d/n^2 \le x \le 1$. 
Since $0\le (P_n^{(\alpha,\beta)})'(t)\le (P_n^{(\alpha,\beta)})'(1)$ for any $t\in[x,1]$,
\[
P_n^{(\alpha,\beta)}(x)
= P_n^{(\alpha,\beta)}(1) - \int_x^1 (P_n^{(\alpha,\beta)})'(t)\, dt 
\ge P_n^{(\alpha,\beta)}(1) - \frac{d}{n^2} (P_n^{(\alpha,\beta)})'(1).
\]
For the reverse inequality we note that $P_n^{(\alpha,\beta)}(x)\le P_n^{(\alpha,\beta)}(1)\approx n^\alpha$.
\end{proof}




Darboux formula allows us to compute the $L_p(\mu_{\alpha,\beta})$-norms of Jacobi polynomials.
Let $\alpha,\beta$ be such that $\min\{\alpha,\beta\}> -1/2$ and put
\[
p(\alpha,\beta)=\frac{4(\gamma+1)}{2\gamma+3}, \quad
q(\alpha,\beta)=\frac{4(\gamma+1)}{2 \gamma +1}, \quad\text{ where }\gamma=\max\{\alpha,\beta\}.
\]
Notice that $p(\alpha,\beta)$ and $q(\alpha,\beta)$ are conjugate exponents.
 We have (see \cite{Leves-Kus}) that, for $n\ge 2$,
\begin{equation}
\label{result:JacobiBasisNorm}
\begin{aligned}
 \Vert p_n^{(\alpha,\beta)} \Vert_{L_p(\mu_{\alpha,\beta})}
&\approx n^{1/2}\Vert P_n^{(\alpha,\beta)} \Vert_{L_p(\mu_{\alpha,\beta})} \\
&\approx
 \begin{cases} 1, & \text{ if } 1\le p<q(\alpha,\beta), \\
 (\log n)^{1/p}, & \text{ if } p=q(\alpha,\beta),\\
 n^{(2 \gamma+1)/2-2( \gamma+1)/p}, & \text{ if } q(\alpha,\beta)<p<\infty.
\end{cases}
\end{aligned}
\end{equation}


\begin{lemma}
\label{result:NormalizationAndBoundedness} 
Let $1\le p< \infty$ and  suppose $\min\{\alpha, \beta\}>-1/2$.  Then
\begin{itemize}
\item[(a)] The Jacobi polynomials of indices $\alpha$ and $\beta$ form a $M$-bounded basis for $L_p(\mu_{\alpha,\beta})$ if and only if
$p(\alpha,\beta)<p<q(\alpha,\beta)$.
\item[(b)] If $p(\alpha,\beta)<p<q(\alpha,\beta)$, then both $(p_n^{(\alpha,\beta)})_{n=0}^\infty$ and $(n^{1/2} P_n^{(\alpha,\beta)})_{n=0}^\infty$
 form a semi-normalized $M$-bounded basis for $L_p(\mu_{\alpha,\beta})$.
\end{itemize}
\end{lemma}

\begin{proof} 
Let $q$ be the conjugate exponent of $p$. Equation \eqref{result:JacobiBasisNorm} yields
\begin{equation*}
\Vert p_n^{(\alpha,\beta)} \Vert_p
\Vert p_n^{(\alpha,\beta)} \Vert_q
\approx
 \begin{cases}
  n^{(2 \gamma+1)/2-2( \gamma+1)/q}, & \text{ if } 1\le p < p(\alpha,\beta),\\
   (\log n)^{1/q}, & \text{ if } p=p(\alpha,\beta),\\
  1, & \text{ if }  p(\alpha,\beta) <p<q(\alpha,\beta), \\
 (\log n)^{1/p}, & \text{ if } p=q(\alpha,\beta),\\
 n^{(2 \gamma+1)/2-2( \gamma+1)/p}, & \text{ if } q(\alpha,\beta)<p.
\end{cases}
\end{equation*}
Appealing to \eqref{EarlyTest}, the proof is over.
\end{proof}

\begin{remark}
Notice that the range of indices $p$ for which the Jacobi polynomials with $\min\{\alpha,\beta\}>-1/2$ form a bounded basis for $L_p(\mu_{\alpha,\beta})$ coincides with the range of indices for which they are a Schauder basis of $L_p(\mu_{\alpha,\beta})$ with the natural order (see ~\cite{Pollard}).
\end{remark}


\begin{remark}
Lemma~\ref{result:NormalizationAndBoundedness}, combined with either Theorem~\ref{NotSimultaneouslyBounded} or Corollary~9 in \cite{KadPel1962}, gives that Jacobi polynomials with $\min\{\alpha,\beta\}>-1/2$ are not an unconditional basis for $L_p(\mu_{\alpha,\beta})$ unless $p=2$.
This result can also be obtained from Proposition~4 in \cite{KonigNielsen1994}.
\end{remark}

Lemma~\ref{result:Random} provides a tool for checking if a given basis is a suitable candidate to be quasi-greedy, and leads us to compare norms of the form
$\Vert \sum_{j\in A} \xx_j \Vert_p$ with norms of the form $\Vert (\sum_{j\in A} |\xx_j|^2)^{1/2} \Vert_p$. 
In this direction, we state the following result.

\begin{proposition}
\label{result:AverageJacobiNorm} 
Let  $\alpha$, $\beta$, and $p$ be such that $\min\{\alpha, \beta\}>-1/2$ and
  $1\le p<q(\alpha,\beta)$.
Then, for $A\subseteq \NN$ finite,
\[
\left\Vert\left(\sum_{n\in A} \left(n^{1/2} P_n^{(\alpha,\beta)}\right)^2\right)^{1/2}\right\Vert_{L_p(\mu_{\alpha,\beta})} \approx 
\left\Vert\left(\sum_{n\in A} \left(p_n^{(\alpha,\beta)}\right)^2\right)^{1/2}\right\Vert_{L_p(\mu_{\alpha,\beta})} \approx |A|^{1/2}.
\]
\end{proposition}

\begin{proof}
Just combine Proposition~\ref{result:LdosSumEstimate} with~\eqref{result:JacobiBasisNorm}.
\end{proof}


Proposition~\ref{result:AverageJacobiNorm} says that the expected value of the $L_p(\mu_{\alpha,\beta})$-norms \[
\left\Vert \sum_{n\in A} \varepsilon_n n^{1/2} P_n^{(\alpha,\beta)}\right\Vert_{L_p(\mu_{\alpha,\beta})},
\]
when $(\varepsilon_n)_{n\in A}$ runs over all possible signs $\{\pm 1\}^A$,
is of the order of $|A|^{1/2}$ (see Theorem 6.2.13 in  \cite{AlbiacKalton2006}). In the next proposition we find norms that deviate significantly from the average value for $p\not= 2$.


\begin{proposition}
\label{result:CCJacobiNorm} 
Let $\alpha$, $\beta$, and $p$ be such $\min\{\alpha, \beta\} >-1/2$ and 
$p(\alpha,\beta)<p<q(\alpha,\beta)$.
Then
\[
\left\Vert\sum_{n=0}^{N-1} (N+2n)^{1/2} P_{N+2n}^{(\alpha,\beta)} 
\right\Vert_{L_p(\mu_{\alpha,\beta})} \approx N^\omega, \quad N\in\NN,
\]
where $\omega=\max\{{(2\alpha+3)/2-2(\alpha+1)/p},{(2\beta+3)/2-2(\beta+1)/p}\}$.
\end{proposition}

\begin{proof}
 Since $P_n^{(\alpha,\beta)}(-x)=(-1)^n P_n^{(\beta,\alpha)}(x)$ it suffices to prove the estimate
\begin{equation*}
\label{eq:esti-alpha}
I_N := \int_0^1\left|\sum_{n\in A_N} n^{1/2} P_{n}^{(\alpha,\beta)}(x)\right|^p \, d\mu_{\alpha,\beta}(x) \approx N^\sigma,
\end{equation*}
where $A_N=\{N+2n : 0\le n\le N-1\}$ and $\sigma=p(2\alpha+3)/2-2(\alpha+1)$. Notice that the hypothesis $p(\alpha,\beta)<p<q(\alpha,\beta)$ implies $0<\sigma<p$.

Consider $d>0$ as in Lemma~\ref{result:NearOneEstimate} and choose $\theta_N\in(0,\pi)$ such that
\[
\cos(\theta_N)=x_N=1-\frac{d}{(3N-2)^2}.
\]
Write $I_N=J_N+K_N$, where
\begin{align*}
J_N &= \int_{x_N}^1 \left|\sum_{n\in A_N} n^{1/2} 
P_{n}^{(\alpha,\beta)}(x)\right|^p \, d\mu_{\alpha,\beta}(x),\quad \text{and}\\
K_N &= \int_{0}^{x_N} \left|\sum_{n\in A_N} n^{1/2} P_{n}^{(\alpha,\beta)}(x) \right|^p \, d\mu_{\alpha,\beta}(x).
\end{align*}
Since $P_n(x)\approx n^\alpha$ for  $n\in A_N$, $x\in[x_N,1]$, and $N\in\NN$,
\[
J_N\approx \left(\sum_{n\in A_N} n^{\alpha+1/2} \right)^p \int_{x_N}^1 (1-x)^\alpha\, dx \approx N^{p(\alpha+3/2)} N^{-2(\alpha +1)} = N^{\sigma}.
\]
Let $k$, $\phi$, and $E_n$ be as is Theorem~\ref{result:Darboux}.
We have $K_N \le 2^{p-1} (L_N+ M_N)$, where
\begin{align*}
L_N &= \int_0^{x_N}\left| \sum_{n\in A_N} E_{n}(\arccos x) \right|^p \, d\mu_{\alpha,\beta}(x),
\quad\text{and} \\
M_N &= \int_0^{x_N}\left| k(\arccos x) \sum_{n\in A_N} \cos(n \arccos x+\phi(\arccos x) )\right|^p \, d\mu_{\alpha,\beta}(x).
\end{align*}

Since $\theta_N\approx N^{-1}$, there exists $\delta>0$ such that $\delta/N\le \theta_N$. By Theorem~\ref{result:Darboux},
\[
|E_n(\theta)|\lesssim \frac{1}{n} (\sin \tfrac{\theta}{2})^{-\alpha-3/2}
\]
for $n\in A_N$, $\theta\in[\theta_N,\pi/2]$, and $N\in\NN$.
The change of variable $x=\cos\theta$ yields
\[
L_N\lesssim\left( \sum_{n\in A_N} \frac{1}{n}\right)^p \int_{\theta_N}^{\pi/2} 
( \sin\tfrac{\theta}{2})^ {-1-\sigma} \, d\theta
\approx \int_{\theta_N}^\infty \theta^{-1-\sigma}\, d\theta 
= \frac{1}{\sigma} \theta_N^{-\sigma} \approx N^\sigma.
\]
Using the change of variable $x=\cos\theta$ and the formula
\[
\left| \sum_{n\in A_N} \cos(n\theta+\phi(\theta))\right|
= \frac{|\sin(N\theta) \cos(2N\theta+\phi(\theta))|}{\sin \theta},
\]
which is obtained by taking into account that we are adding the real part of a geometric sum, gives
\begin{align*}
M_N &\approx \int_{0}^{\pi/2} |\sin(N\theta) \cos(2N\theta+\phi(\theta))|^p 
(\sin\tfrac{\theta}{2})^{-1-\sigma}\,d\theta \\
&\approx \int_{0}^{\pi/2} |\sin(N\theta) \cos(2N\theta+\phi(\theta))|^p \theta^{-1-\sigma}\,d\theta \\
&= N^\sigma \int_{0}^{N\pi/2} |\sin(u) \cos(2u+\phi(u/N))|^p u^{-1-\sigma} \, du \\
& \approx N^\sigma.
\end{align*}
To deduce the last step in the estimate of $M_{N}$ we used that
\[
 \int_0^\infty |\sin(u)|^p  u^{-1-\sigma} \, du<\infty
\]
and the Dominated Convergence theorem.
\end{proof}

We are now in a position to complete the proof Theorem~\ref{result:MainTheorem} as advertised.

\begin{proof}[Proof of Theorem~\ref{result:MainTheorem}]
Assume that the $L_p(\mu_{\alpha,\beta})$-normalized sequence of Jacobi polynomials 
is a quasi-greedy basis for $L_p(\mu_{\alpha,\beta})$.
Then, thanks to Lemma~\ref{result:NormalizationAndBoundedness}(a), we have $p(\alpha,\beta)<p<q(\alpha,\beta)$.

By Theorem~\ref{result:Perturbing} and Lemma~\ref{result:NormalizationAndBoundedness}(b), the basis
$(n^{1/2} P_n^{(\alpha,\beta)})_{n=0}^\infty$ is also quasi-greedy for $L_p(\mu_{\alpha,\beta})$. 
Combining Proposition~\ref{result:Democracy} 
(or Lemma~\ref{result:Random} together with Proposition~\ref{result:AverageJacobiNorm}) with
Proposition~\ref{result:CCJacobiNorm} we obtain
$N^\omega \approx N^{1/2}$ for $N\in\NN$.
Therefore, $\omega=1/2$, i.e., $p=2$.
\end{proof}

% -----------------------
\subsection*{Acknowledgements}
The first two authors were partially supported by the Spanish Research Grant \textit{An\'alisis Vectorial, Multilineal y Aplicaciones}, reference number MTM2014-53009-P, and the last two authors were partially supported by the Spanish Research Grant \textit{Ortogonalidad, Teor\'ia de la Aproximaci\'on y Aplicaciones en F\'isica Matem\'atica}, reference number MTM2015-65888-C4-4-P. The first-named author also acknowledges the support of Spanish Research Grant \textit{Operators, lattices, and structure of Banach spaces}, 
with reference  MTM2016-76808-P.
% -----------------------

%%%%%%%%%%%%%%%%%%%%%%%%%%%%%%%%%
% Bibliography
%%%%%%%%%%%%%%%%%%%%%%%%%%%%%%%%%

% -----------------------
\begin{thebibliography}{99}
% \parskip=0.7pt

\newcommand{\RMIauthor}[1]{\textsc{#1}:}
\newcommand{\RMIpaper}[1]{\textnormal{#1},}
\newcommand{\RMIjournal}[1]{\textit{#1}}
\newcommand{\RMIbook}[1]{\textit{#1},}

\bibitem{AA2015}
\RMIauthor{Albiac, F. and Ansorena, J.L.}
\RMIpaper{Lorentz spaces and embeddings induced by almost greedy bases in Banach spaces}
\RMIjournal{Constr. Approx.} 
\textbf{43} (2016), 197--215.

\bibitem{AlbiacKalton2006}
\RMIauthor{Albiac, F. and Kalton, N.J.}
\RMIbook{Topics in Banach space theory, Second Edition}
Graduate Texts in Mathematics, 233, Springer, 2016.

\bibitem{BePe1958}
\RMIauthor{Bessaga, C. and Pe{\l}czy{\'n}ski, A.}
\RMIpaper{On bases and unconditional convergence of series in Banach spaces}
\RMIjournal{Studia Math.} 
\textbf{17} (1958), 151--164.

\bibitem{CoFe1998}
\RMIauthor{C{\'o}rdoba, A. and Fern{\'a}ndez, P.}
\RMIpaper{Convergence and divergence of decreasing rearranged Fourier series}
\RMIjournal{SIAM J. Math. Anal.} 
\textbf{29} (1998), no. 5, 1129--1139.

\bibitem{DKKT2003}
\RMIauthor{Dilworth, S.J., Kalton, N.J., Kutzarova, D. and Temlyakov, V.N.}
\RMIpaper{The thresholding greedy algorithm, greedy bases, and duality}
\RMIjournal{Constr. Approx.} 
\textbf{19} (2003), 75--597.

\bibitem{DST2012}
\RMIauthor{Dilworth, S.J., Soto-Bajo, M. and Temlyakov, V.N.}
\RMIpaper{Quasi-greedy bases and Lebesgue-type inequalities}
\RMIjournal{Studia Math.} 
\textbf{211} (2012), 41--69.

\bibitem{Gaposkin1958}
\RMIauthor{Gapo{\v{s}}kin, V.F.}
\RMIpaper{On unconditional bases in $L^{p}$ ($p>1$) spaces (Russian)}
\RMIjournal{Uspehi Mat. Nauk} 
\textbf{13} (1958), no. 4(82), 179--184.

\bibitem{JohnsonLindes2001}
\RMIauthor{Johnson, W.B. and Lindenstrauss, J.}
\RMIpaper{Basic concepts in the geometry of Banach spaces}
In \RMIbook{Handbook of the geometry of Banach spaces, Vol. I, {\rm 1--84}}
North-Holland, Amsterdam, 2001.

\bibitem{KadPel1962}
\RMIauthor{Kadec, M.I. and Pe{\l}czy{\'n}ski, A.}
\RMIpaper{Bases, lacunary sequences and complemented subspaces in the spaces $L_{p}$}
\RMIjournal{Studia Math.} 
\textbf{21} (1961/1962), 161--176.

\bibitem{KonigNielsen1994}
\RMIauthor{K{\"o}nig, H. and Nielsen, N.J.}
\RMIpaper{Vector-valued $L_p$-convergence of orthogonal series and Lagrange interpolation}
\RMIjournal{Forum Math.} 
\textbf{6} (1994), no. 2, 183--207.

\bibitem{KonyaginTemlyakov1999}
\RMIauthor{Konyagin, S.V. and Temlyakov, V.N.}
\RMIpaper{A remark on greedy approximation in Banach spaces}
\RMIjournal{East J. Approx.} 
\textbf{5} (1999), no. 3, 365--379.

\bibitem{Leves-Kus}
\RMIauthor{Levesley, S. and Kushpel, A.K.}
\RMIpaper{On the norm of the Fourier-Jacobi projection}
\RMIjournal{Numer. Funct. Anal. Optim.} 
\textbf{22} (2001), 941--952.


\bibitem{LinRos1969}
\RMIauthor{Lindenstrauss, J. and Rosenthal, H.P.}
\RMIpaper{The $\LL_p$ spaces}
\RMIjournal{Israel J. Math.} 
\textbf{7} (1969), 325--349.

\bibitem{HMVZ2008}
\RMIauthor{H{\'a}jek, P., Montesinos Santaluc{\'{\i}}a, V., Vanderwerff, J. and Zizler, V.}
\RMIbook{Biorthogonal systems in Banach spaces}
CMS Books in Mathematics/Ouvrages de Math\'ematiques de la SMC, 26, Springer, New York, 2008.

\bibitem{Nielsen2007}
\RMIauthor{Nielsen, M.}
\RMIpaper{An example of an almost greedy uniformly bounded orthonormal basis for $L_p(0,1)$}
\RMIjournal{J. Approx. Theory} 
\textbf{149} (2007), 188--192.

\bibitem{Pollard}
\RMIauthor{Pollard, H.}
\RMIpaper{The mean convergence of orthogonal series III}
\RMIjournal{Duke Math. J.} 
\textbf{16} (1949), 189--191.

\bibitem{Szego}
\RMIauthor{Szeg\H{o}, G.}
\RMIbook{Orthogonal polynomials (3rd edition)}
Amer. Math. Soc. Colloq. Publ., vol. 23, Amer. Math. Soc., Providence, R.I., 1967.

\bibitem{Telmyakov1998}
\RMIauthor{Temlyakov, V.N.}
\RMIpaper{Greedy algorithm and $m$-term trigonometric approximation}
\RMIjournal{Constr. Approx.} 
\textbf{14} (1998), 569--587.

\bibitem{Wo2000}
\RMIauthor{Wojtaszczyk, P.}
\RMIpaper{Greedy algorithm for general biorthogonal systems}
\RMIjournal{J. Approx. Theory} 
\textbf{107} (2000), 293--314.

\bibitem{Zippin1966}
\RMIauthor{Zippin, M.}
\RMIpaper{On perfectly homogeneous bases in Banach spaces}
\RMIjournal{Israel J. Math.} 
\textbf{4} (1966), 265--272.

\end{thebibliography}
% -----------------------

% ------------------------------------------------------------------------
\end{document}
% ------------------------------------------------------------------------

